% !PRINTED 278
\weberpage{278}
Under this assumption the formulae of \S.\,57 lead to the setting up
of an octahedral group.

The generating substitutions are:
\[
\Theta = \begin{pmatrix} \gamma^{\frac{p^2-1}{8}}, & 0 \\ 0, & \gamma^{-\frac{p^2-1}{8}} \end{pmatrix}, \quad
\psi = \begin{pmatrix} 0, & \gamma^{\frac{p^2-1}{4}} \\ \gamma^{\frac{p^2-1}{4}}, & 0 \end{pmatrix}
\]
\weq{8}{
\chi = \begin{pmatrix} \dfrac{1}{\sqrt{2}}\,\gamma^{-\frac{p^2-1}{8}}, & \dfrac{1}{\sqrt{2}}\,\gamma^{-\frac{p^2-1}{8}} \\[2mm] \dfrac{-1}{\sqrt{2}}\,\gamma^{\frac{p^2-1}{8}}, & \dfrac{1}{\sqrt{2}}\,\gamma^{\frac{p^2-1}{8}} \end{pmatrix}.
}
This group is real when $p\equiv 1\ (\mathrm{mod}\ 8)$, and imaginary, but contained in $\Gamma_p$, when $p\equiv -1\ (\mathrm{mod}\ 8)$.

The first example is $p=7$. Here $-1$ is a quadratic non-residue, and one can therefore set $\varepsilon=\sqrt{-1}=i$. If we take $\gamma=2+i$, $\gamma^7=2-i$, $\gamma^6=2\,(1+i)$, $\gamma^8=5$, then $\gamma$ is a primitive root, since $\gamma,\gamma^2,\dots,\gamma^7$ are different from one another and $5$ is a real primitive root of $7$, so that every non-vanishing number of the field $\mathfrak{E}_{2,7}$ can be represented in the form
\[
5^{\mu}\gamma^{\nu} = \gamma^{8\mu+\nu}
\]
if
\[
\begin{gathered}
\mu = 0,\ 1,\ 2,\ 3,\ 4,\ 5,\\
\nu = 0,\ 1,\ 2,\ 3,\ 4,\ 5,\ 6,\ 7
\end{gathered}
\]
is set.

Further, we then have $\sqrt{2}=3$, and consequently the substitutions $\Theta,\psi,\chi$ are:
\weq{9}{
\begin{pmatrix} 2\,(1+i), & 0 \\ 0, & 2\,(1-i) \end{pmatrix}, \
\begin{pmatrix} 0, & i \\ i, & 0 \end{pmatrix}, \
\begin{pmatrix} 4\,(1-i), & 4\,(1-i) \\ -4\,(1+i), & 4\,(1+i) \end{pmatrix}.
}
If one wishes to pass to the real form, one must proceed according to \S.\,69. One passes from the substitutions $B$ to the $A$ by setting $\nu=\gamma^3=2-3i$, and obtains from (9) for $\Theta,\psi,\chi$ in the group $\Lambda_p$:
\weq{10}{
\begin{pmatrix} 2+2i, & 0 \\ 0, & 2-2i \end{pmatrix}, \
\begin{pmatrix} 0, & 3-2i \\ 3+2i, & 0 \end{pmatrix}, \
\begin{pmatrix} 4-4i, & 1-4i \\ 1+4i, & 4+4i \end{pmatrix}.
}
Here, further,
\[
R = \begin{pmatrix} 4, & 4i \\ i, & 1 \end{pmatrix}
\]
%CONT>>
